%postscript
\documentclass{article}
%\usepackage{fancybox}% http://ctan.org/pkg/fancybox
\usepackage{pgfornament}
\usetikzlibrary{calc}
\usepackage{epsfig}
\usepackage{pst-node}% http://ctan.org/pkg/pstricks
\usepackage{xcolor}% http://ctan.org/pkg/xcolor
\usepackage{eso-pic}% http://ctan.org/pkg/eso-pic
\usepackage[T1,T2A]{fontenc}         % внутрішнє кодування шрифтів (може бути декілька);
                                          % вказане останнім діє по замовчуванню;
                                          % кириличне має співпадати з заданим в ukrhyph.tex.
\usepackage[utf8]{inputenc}       % кодування документа; замість cp866nav
                                          % може бути cp1251, koi8-u, macukr, iso88595, utf8.
\usepackage[english,ukrainian]{babel} % національна локалізація; може бути декілька
\begin{document} 
\pagestyle{empty}
\parindent 0pt

\begin{titlepage}
  \AddToShipoutPictureBG*{%
    \begin{pspicture}(\paperwidth,\paperheight)
      \psset{xunit=\the\paperwidth,yunit=\the\paperheight,nodesep=2cm}
      \pnode(1,1){page.northeast}%
      \pnode(0,0){page.southwest}%

      \psframe[linestyle=dashed,linewidth=10pt,linecolor=orange!50]([angle=45]page.southwest)([angle=-135]page.northeast)
    \end{pspicture}
  }

\newcommand{\hl}{\bigskip\hrule\bigskip}
\pagestyle{empty}

\setlength{\unitlength}{\textwidth}

\begin{picture}(0,0)(0,-0.05)
\put(-0.09,0){\epsfig{file=gerb145.eps,height=2cm}}
\put(0.5,0.04){\makebox(0,0){\Large http://www.lic145.kiev.ua/}}
\put(0.5,0.08){\makebox(0,0){\Large Київський природничо-науковий ліцей №145}}
\put(0.5,-0.01){\makebox(0,0){\rule{\textwidth}{0.5pt}}}
%\put(1,0){\makebox(0,0)[rb]{\epsfig{file=pics/gerb145.eps,height=1.8cm}}}
\end{picture}

\begin{center}
{\huge \bf Переваги і недоліки навчання за кордоном\\
для випускників ліцею №145}
\bigskip

{\Large 
{\bfseries {\it Олексій Леонтьєв}} \\
Tokyo University\\Graduate School of Mathematical Sciences, M.A.}
\end{center}

{\small Зміст: Протягом 8 років десь близько 40 випускників 145 ліцею стали студентами Тайванського National
Chiao Tung University. Їх долі склалися по-різному: частина з них продовжує навчання в престижних університетах на Тайвані
і не тільки (NTU (Тайвань), Chalmers U (Швеція), Princeton (США) тощо), дехто зараз працює в галузях торгівлі, електронної
інженерії, програмування тощо. На прикладі мене і моїх знайомих я спробую пояснити переваги і недоліки навчання за кордоном (більш
точно, переваги і недоліки вступу в NCTU), на що треба очікувати потрапивши туди і які перспективи відкриває цей крок.

В 2009 році я потрапив на Тайвань, будучи випускником ліцею №145 разом із п’ятьма іншими хлопцями і дівчатами. Після цього я навчався
в NCTU протягом 3.5 років (отримавши звання бакалавра з двох факультетів: програмування і прикладної математики)
, провів семестр як студент за обміном в HKU (Гонконг), а також два роки в Tokyo University (Японія) де
отримав звання магістра в області теорії представлень груп Лі. Володію англійською, китайською і японською мовами.
}\\

\bigskip
\begin{center}
\begin{tikzpicture}[every node/.style={inner sep=0pt}]
\node[text width=8cm,align=center](Text){
\begin{tabular}{ll}
День:&(empty)\\
Час:&(empty)\\
Місце:&(empty)\\
Мова:&російська\\
\end{tabular}
};
\node[shift={(-1cm,1cm)},anchor=north west](CNW)
at (Text.north west) {\pgfornament[width=1.75cm]{61}};
\node[shift={(1cm,1cm)},anchor=north east](CNE)
at (Text.north east) {\pgfornament[width=1.75cm,symmetry=v]{61}};
\node[shift={(-1cm,-1cm)},anchor=south west](CSW)
at (Text.south west) {\pgfornament[width=1.75cm,symmetry=h]{61}};
\node[shift={(1cm,-1cm)},anchor=south east](CSE)
at (Text.south east) {\pgfornament[width=1.75cm,symmetry=c]{61}};
\pgfornamenthline{CNW}{CNE}{north}{82}
\pgfornamenthline{CSW}{CSE}{south}{82}
%%\pgfornamentvline{CNW}{CSW}{west}{82}
%%\pgfornamentvline{CNE}{CSE}{east}{82}
\end{tikzpicture}
\end{center}
\begin{center}
\bigskip
Запрошуємо усіх бажаючих!
\end{center}
\end{titlepage}
\end{document}
